\section{Common notation and the native target}\label{sec:model}

\paragraph{Matrices and information.}
The dimension is $d$, the number of bits in an entry is $b$, and
$P=d^2b$. An input is a bit matrix $x=(x_{ijt})$ with $i,j<d$ and $t<b$.
Its row-major encoding writes bit $(i,j,t)$ at position $(di+j)b+t$;
the required output puts it at $(dj+i)b+t$. We use base-two Shannon entropy.
For a random input $X$, its deficit is $D_X=P-H(X)$. If $X$ is uniform on a
nonempty set $F$, this is $D=P-\log_2|F|$. A helper $Z$ is a deterministic
function of the whole input, supplied to the participants specified in a
particular model. If $X$ is supported on $F$, the bound $|Z(F)|\le2^g$
implies $H(Z)\le g$. Each theorem specifies access to the helper; it does
not supply a tape implementation of that access.

The finite coding results allow arbitrary $b$ unless stated otherwise.
For arbitrary input laws they use $D_X$; the cardinality formula for $D$
requires the uniform law on $F$. All uses of $o(P)$ are along the indicated
size sequence.

\paragraph{Native machines.}
We use the unchanged model of~\cite{FirstPaper2026}. A machine has $q+1$
states, $k$ two-way infinite work tapes over $g+1$ symbols (one blank), a
read-only two-way input tape, and a one-way write-only output tape. In one
transition each input or work head moves by $-1$, $0$ or $1$; the machine may append at most
one output bit. The input lies in cells $0,\ldots,P-1$, with an end marker
elsewhere. All heads start at zero, work tapes are blank, and the output is
empty. The transition table, state set, alphabet and tape count are finite
and independent of the input for a uniform machine. Write
$G=\max(1,\ceil{\log_2(g+1)})$ for the charge in bits per work cell.
``Computes within $t$'' means that an actual halting configuration with exactly
the specified output is reached after at most $t$ transitions.
\upformal{Transpose.Statement}{Transpose.TM,Transpose.ComputesWithin}
\upformal{Transpose.Tape.DefsM}{Transpose.Tape.cellBits}

\begin{definition}[The unchanged decision target]\label{def:target}
Let
\[
 b(d)=\max(1,\ceil{\log_2d}),\qquad P(d)=d^2b(d),\qquad d\ge1.
\]
Statement \textup{(T)} asserts that there exist one native machine $M$ and one
constant $c\in\N$ such that, for every $d\ge1$ and every input matrix,
$M$ halts with its raw transpose within $c(P(d)+1)$ steps.
\end{definition}
\upformal{Transpose.Statement}{Transpose.entryBits,Transpose.inputLength,Transpose.encode,Transpose.transpose,Transpose.TransposesWithin,Transpose.LinearTimeTransposable}

The Lean files use $n$ for the dimension in many declarations; the manuscript
consistently writes $d$. Formally, the width is
$\operatorname{Nat.log2}(d-1)+1$, including the convention at $d=1$.
The original target and checker are reused without alteration;
\textup{(T)} remains open.

\paragraph{The multiplication connection (cited background).}
Write $\lg x=\max(1,\ceil{\log_2x})$ and let $\lg^*x$ be the number of
iterations of this rounded logarithm needed to reach $1$.
Harvey and van der Hoeven~\cite[Corollary~6.9]{HvdH2025} reduce transposition
with entry width $b$ to multiplication with an iterated-logarithm factor
$\ell_0=\max(1,\lg^*d-\lg^*b)$. At $b=\lg d=b(d)$ this factor is one.
The encoding conventions are reconciled in the earlier paper's discussion
of its Proposition on $(T)$ and $1$-logarithmic transposition.
Consequently a negative resolution of \textup{(T)} would exclude linear-time
multiplication in the multitape model. A positive resolution of \textup{(T)}
would not by this implication produce a multiplier. For binary entries
$b=1$, the cited reduction from a linear multiplier instead gives an
$O(P\log^*P)$ transposer.

\paragraph{Scope of the sequence statements.}
Native exclusion theorems below can allow a machine chosen separately for each
size, while fixing its number of work tapes and cell-bit charge and bounding
the logarithm of its state count by $o(P)$. Such statements are stronger than
their fixed-machine special cases, but do not allow input-dependent advice
without charging it. The correctness set may have deficit $o(P)$; geometric
hypotheses refer to actual runs on that set. The required universal and
subsequential quantifiers are given explicitly in the motion section.
